% Part: lambda-calculus
% Chapter: lambda-definability
% Section: partial-recursive-functions

\documentclass[../../../include/open-logic-section]{subfiles}

\begin{document}

\olfileid{lam}{ldf}{par}
\olsection{部分递归函数是\usetoken{S}{lambda definable}}

部分递归函数是指从基本函数出发，通过复合、原始递归和无界极小化得到的函数。
它们与一般递归函数的区别在于，无界搜索中使用的函数不要求是正则的。
不要求正则性意味着由极小化定义的函数可能无定义。

乍看之下，似乎可以用证明（全的）一般递归函数都是 $\lambd$-可定义的同样方法，
来证明所有部分递归函数都是 $\lambd$-可定义的。
例如，若 $f$ 和 $g$ 分别由项 $F$ 和~$G$ $\lambd$-定义，则 $f$ 与 $g$ 的复合
可由 $\lambd[x][F (G x)]$ $\lambd$-定义。
然而，当函数为部分函数时，便成了问题。
当 $g(x)$ 无定义时，即 $G x$ 没有范式；
在大多数情况下，这意味着 $F (G x)$ 也没有范式，而这正是我们所期望的。
但考虑 $F$ 为 $\lambd[x][\lambd[y][y]]$ 的情形，此时 $F (G x)$ 却有范式（$\lambd[y][y]$）。

这个问题并非无法克服，存在一些方法可以 $\lambd$-定义所有部分递归函数，
使得无定义的值由无范式的项来表示。
不过，这些方法比我们处理一般递归函数所采用的方法更为复杂，也不那么直观。
我们在此不加证明地陈述以下定理：

\begin{thm}
  所有部分递归函数都是 $\lambd$-可定义的。
\end{thm}

\end{document}